\section{双曲线}

\subsection{双曲线的定义与标准方程}

\begin{definition}
    平面上两点 $F_1, F_2$ 的距离之差为常数的点的轨迹称为双曲线。
\end{definition}

\begin{theorem}
    双曲线的标准方程为 $\frac{x^2}{a^2} - \frac{y^2}{b^2} = 1 (a>0, b>0)$ 或 $\frac{y^2}{a^2} - \frac{x^2}{b^2} = 1 (a>0, b>0)$。
\end{theorem}

\begin{example}
    已知曲线 $C: m x^2+n y^2=1$ ，下列结论正确的是 （\qquad）
    \begin{itemize}
        \item 若 $m>n>0$ ，则 $C$ 是椭圆，其焦点在 $y$ 轴上
        \item 若 $m>0>n$ ，则 $C$ 是双曲线，其焦点在 $x$ 轴上
        \item 若 $m=n$ ，则 $C$ 是圆
        \item 若 $m=0, n \neq 0$ ，则 $C$ 是两条直线
    \end{itemize}
\end{example}

\vspace{10em}

\subsection{含有参数的双曲线判断}

\begin{example}
    ＂ $1<k<5$＂是方程＂$\frac{x^2}{k-1}+\frac{y^2}{5-k}=1$ 表示椭圆＂的 （\qquad）
    \begin{tasks}(4)
        \task 充分不必要条件
        \task 必要不充分条件
        \task 充要条件
        \task 既不充分又不必要条件
    \end{tasks}
\end{example}

\vspace{10em}

\begin{example}
    当 $m<2$ 时，方程 $\frac{x^2}{5-m}+\frac{y^2}{2-m}=1$ 表示的曲线是（\qquad）
    \begin{tasks}(4)
        \task 焦点在 $x$ 轴的椭圆
        \task 焦点在 $y$ 轴的椭圆
        \task 双曲线
        \task 圆
    \end{tasks}
\end{example}

\vspace{10em}

\begin{example}
    曲线 $C: \frac{x^2}{m^2-1}-\frac{y^2}{m}=1$ 是焦点在 $x$ 轴上的双曲线，则 $m$ 的取值范围为 \underline{\hspace{3cm}}。
\end{example}

\vspace{10em}

\begin{example}
    已知方程 $\frac{m x^2}{2-m}+\frac{y^2}{m-3}=1$ 所表示的曲线为双曲线，则 $m$ 的取值范围为 \underline{\hspace{3cm}}。
\end{example}

\vspace{10em}

\begin{example}
    若双曲线 $\frac{x^2}{3-k}+\frac{y^2}{k-1}=1$ 的焦点在 $y$ 轴上，则实数 $k$ 的取值范围为 \underline{\hspace{3cm}}。
\end{example}

\vspace{10em}

\begin{example}
    若方程 $m y^2+(2+m) x^2=1$ 表示焦点在 $x$ 轴上的双曲线，则实数 $m$ 的取值范围为 \underline{\hspace{3cm}}。
\end{example}

\vspace{10em}

\begin{theorem}
    \begin{tabular}{| p{2cm} | >{\centering\arraybackslash}p{8cm} | >{\centering\arraybackslash}p{6cm} |}
        \hline
        \textbf{名称}            & \textbf{焦点在 $x$ 轴上}                                                                                                   & \textbf{焦点在 $y$ 轴上}                                \\
        \hline
        第一定义                   & \multicolumn{2}{c|}{平面内一动点 $P$ 与两定点 $F_1, F_2$ 距离之差为常数 (大于 $|F_1F_2|$) 的轨迹}                                                                                                \\
        \hline
        第二定义                   & \multicolumn{2}{c|}{平面内一动点到定点与到准线的距离比是常数的点的轨迹 $e = \frac{MF}{d_1}$}                                                                                                        \\
        \hline
        焦点                     & \multicolumn{2}{c|}{焦点在 $x$ 轴上}                                                                                                                                            \\
        \hline
        标准方程                   & $\frac{x^2}{a^2} - \frac{y^2}{b^2} = 1 (a>0, b>0)$                                                                    & $\frac{y^2}{a^2} - \frac{x^2}{b^2} = 1 (a>0, b>0)$ \\
        \hline
        范围                     & $x \le -a \text{ 或 } x \ge a, y \in \mathbb{R}$                                                                       & $y \le -a \text{ 或 } y \ge a, x \in \mathbb{R}$    \\
        \hline
        顶点                     & $A_1(-a,0), A_2(a,0)$                                                                                                 & $A_1(0,-a), A_2(0,a)$                              \\
        \hline
        轴长                     & \multicolumn{2}{c|}{虚轴长 $=2b$, 实轴长 $=2a$, 焦距 $=|F_1F_2|=2c$, $c^2=a^2+b^2$}                                                                                                \\
        \hline
        焦点                     & $F_1(-c,0), F_2(c,0)$                                                                                                 & $F_1(0,-c), F_2(0,c)$                              \\
        \hline
        焦半径                    & \multicolumn{2}{c|}{$|PF| = |a + ex_0|$}                                                                                                                                   \\
        \hline
        离心率                    & \multicolumn{2}{c|}{$e = \frac{c}{a} = \sqrt{1 + \frac{b^2}{a^2}} (e>1)$}                                                                                                  \\
        \hline
        准线方程                   & $x = \pm \frac{a^2}{c}$                                                                                               & $y = \pm \frac{a^2}{c}$                            \\
        \hline
        渐近线                    & $y = \pm \frac{b}{a}x$                                                                                                & $y = \pm \frac{a}{b}x$                             \\
        \hline
        切线方程                   & $\frac{x_0x}{a^2} - \frac{y_0y}{b^2} = 1$                                                                             & $\frac{y_0y}{a^2} - \frac{x_0x}{b^2} = 1$          \\
        \hline
        通径                     & \multicolumn{2}{c|}{过双曲线焦点且垂直于对称轴的弦长 $|AB| = \frac{2b^2}{a}$ (最短焦弦)}                                                                                                       \\
        \hline
        \multirow{4}{*}{焦点三角形} & \multicolumn{2}{c|}{(1) 由定义可知: $||PF_1| - |PF_2|| = 2a$}                                                                                                                   \\
        \cline{2-3}
                               & \multicolumn{2}{c|}{(2) 焦点直角三角形的个数为八个, 顶角为直角与直角为直角各四个;}                                                                                                                    \\
        \cline{2-3}
                               & \multicolumn{2}{c|}{(3) 焦点三角形面积: $S_{\triangle PFF_1} = b^2 + \tan \frac{\theta}{2} = e \cdot |y|$}                                                                        \\
        \cline{2-3}
                               & \multicolumn{2}{c|}{(4) 离心率: $e = \frac{|F_1F_2|}{||PF_1| - |PF_2||} = \frac{\sin \theta}{\sin \alpha - \sin \beta}$}                                                      \\
        \hline
    \end{tabular}
\end{theorem}